\chapter{Django Admin}
\label{cap:django-admin}

\begin{procedimiento}{Objetivo del capitulo}{Superusuario Django Admin}
Explicar como usar Django Admin para revision y soporte, diferenciandolo del panel operativo del torneo y evitando cambios directos innecesarios.
\end{procedimiento}

Django Admin es una superficie administrativa avanzada. Permite ver y editar modelos registrados de participantes, torneo y comunicaciones. Debe usarse para mantenimiento, revision y contingencias, no como reemplazo del panel de organizacion.

\section{Panel operativo frente a Django Admin}

\begin{tabularx}{\linewidth}{>{\bfseries}p{3.8cm}X}
\toprule
Superficie & Uso recomendado\\
\midrule
Panel de organizacion & Operar el torneo: crear tablero, revisar divisiones, guardar resultados, avanzar fases, gestionar comunicaciones y consultar historial con controles preparados para el evento.\\
Django Admin & Revisar registros, buscar datos, corregir informacion puntual, ejecutar acciones administrativas y auditar modelos cuando el panel operativo no basta.\\
\bottomrule
\end{tabularx}

\begin{operacionsensible}
Django Admin permite cambios directos que pueden saltarse validaciones de la interfaz operativa. Use una cuenta autorizada, trabaje con datos demo cuando este documentando y no modifique registros reales sin respaldo y autorizacion.
\end{operacionsensible}

\section{Acceso e indice de modelos}

\figuraManual{capturas/finales/MU-040-admin-modelos.png}{Indice de Django Admin con aplicaciones y modelos disponibles para administracion.}{fig:bloque4-mu-040}

\begin{procedimiento}{Entrar a Django Admin}{Superusuario Django Admin}
Use este procedimiento solo con una cuenta autorizada para administracion.
\end{procedimiento}

\begin{precondiciones}
\begin{itemize}
    \item El usuario tiene permisos de staff o superusuario segun la politica del proyecto.
    \item La sesion corresponde al entorno correcto.
    \item No se capturan ni comparten credenciales.
\end{itemize}
\end{precondiciones}

\paso{1}{Abra la ruta \ruta{/admin/}.}
\paso{2}{Inicie sesion si Django solicita autenticacion.}
\paso{3}{Revise el indice por aplicaciones.}
\paso{4}{Identifique el modelo antes de abrirlo. No use busquedas al azar sobre datos sensibles.}
\paso{5}{Abra solo el modelo que necesita revisar o corregir.}

\begin{resultadoesperado}
El superusuario ve el indice de modelos registrados y puede entrar a la lista correspondiente sin alterar datos.
\end{resultadoesperado}

\section{Modelos principales}

\begin{tabularx}{\linewidth}{>{\bfseries}p{4cm}X}
\toprule
Area & Modelos y uso habitual\\
\midrule
Participantes & Instituciones, registros escolares, registros universitarios, equipos e integrantes. Se usan para revisar datos de inscripcion, asistencia y sincronizacion al cuadro oficial.\\
Torneo & Ediciones, reglas, fases, competencias por division, grupos, entradas de grupo, batallas, entradas de batalla, estados de equipo e historial. Se usan para soporte y auditoria del cuadro.\\
Comunicaciones & Plantillas, destinatarios, lotes y logs. Se usan para auditar plantillas, destinatarios y resultados de simulaciones, pruebas o envios.\\
\bottomrule
\end{tabularx}

\section{Registros de participantes}

\figuraManual{capturas/finales/MU-041-admin-registros.png}{Lista administrativa de registros con filtros, busqueda y acciones de asistencia.}{fig:bloque4-mu-041}

\begin{procedimiento}{Revisar registros de participantes}{Superusuario Django Admin}
Use este procedimiento para ubicar registros escolares o universitarios y verificar su estado.
\end{procedimiento}

\paso{1}{En el indice, abra \campointerfaz{School registrations} o \campointerfaz{University registrations}, segun la categoria.}
\paso{2}{Use filtros por \campointerfaz{status} y \campointerfaz{edition}.}
\paso{3}{Use busqueda por institucion, robot, responsable o correo cuando necesite ubicar un registro puntual.}
\paso{4}{Revise columnas como institucion, robot, responsable, edicion, estado, solicitud de asistencia enviada y asistencia confirmada.}
\paso{5}{Abra el registro solo si necesita revisar detalle o participantes relacionados.}
\paso{6}{Si debe modificar un dato, cambie solo el campo autorizado y guarde una nota externa de la razon operativa.}

\begin{resultadoesperado}
El superusuario ubica el registro correcto y comprende si esta enviado, confirmado o pendiente, sin alterar otros registros.
\end{resultadoesperado}

\section{Acciones administrativas de asistencia}

Los registros escolares y universitarios incluyen acciones masivas observadas en la interfaz administrativa.

\begin{tabularx}{\linewidth}{>{\bfseries}p{4.6cm}X}
\toprule
Accion & Uso y resultado esperado\\
\midrule
\campointerfaz{Enviar correo de confirmacion de asistencia} & Envia solicitudes a registros seleccionados que no han confirmado asistencia. El mensaje administrativo reporta \campointerfaz{Correos de confirmacion enviados: n. Registros omitidos o con error: m.}\\
\campointerfaz{Sincronizar confirmados al cuadro oficial} & Crea o actualiza equipos a partir de registros con asistencia confirmada. El mensaje administrativo reporta \campointerfaz{Registros sincronizados a equipos: n. Registros omitidos: m.}\\
\bottomrule
\end{tabularx}

\begin{procedimiento}{Usar una accion masiva de registros}{Superusuario Django Admin}
Use una accion masiva solo cuando la lista filtrada contiene exactamente los registros que quiere procesar.
\end{procedimiento}

\paso{1}{Filtre por edicion y estado antes de seleccionar registros.}
\paso{2}{Seleccione manualmente los registros que corresponden.}
\paso{3}{Elija la accion en el selector superior.}
\paso{4}{Revise de nuevo la cantidad seleccionada antes de ejecutar.}
\paso{5}{Ejecute la accion.}
\paso{6}{Lea el mensaje administrativo con procesados y omitidos.}
\paso{7}{Verifique una muestra de registros o equipos resultantes.}

\begin{advertencia}
Las acciones masivas pueden enviar correos o crear equipos desde varios registros. No use \campointerfaz{Seleccionar todos} sin filtrar por edicion y sin validar que no haya registros de prueba mezclados con registros reales.
\end{advertencia}

\section{Administracion del torneo}

\figuraManual{capturas/finales/MU-042-admin-torneo.png}{Lista administrativa de competencias de torneo con filtros por division, estado y edicion.}{fig:bloque4-mu-042}

\begin{procedimiento}{Revisar datos de torneo desde Admin}{Superusuario Django Admin}
Use esta revision cuando necesite comprobar estructura interna del cuadro o buscar una inconsistencia.
\end{procedimiento}

\paso{1}{Abra el modelo de torneo relacionado con la duda: edicion, competencia por division, grupo, batalla, estado de equipo o historial.}
\paso{2}{Use filtros por edicion, division, fase o estado.}
\paso{3}{Use busqueda por nombre de fase, nombre de batalla o robot si el modelo lo permite.}
\paso{4}{Revise relaciones antes de editar: una batalla pertenece a una competencia, una entrada pertenece a una batalla o grupo, y un estado de equipo resume la ubicacion actual.}
\paso{5}{Prefiera corregir resultados desde el panel operativo cuando exista el control correspondiente.}
\paso{6}{Si edita en Admin, revise luego el panel de organizacion para confirmar que la interfaz interpreta correctamente el cambio.}

\begin{fallas}
Si un registro no aparece en la lista esperada, revise primero filtros activos, edicion y division. No cree duplicados para resolver una busqueda fallida.
\end{fallas}

\section{Administracion de comunicaciones}

\begin{procedimiento}{Auditar comunicaciones desde Admin}{Superusuario Django Admin}
Use Django Admin para revisar plantillas, destinatarios, lotes y logs cuando el historial operativo no sea suficiente.
\end{procedimiento}

\paso{1}{Abra \campointerfaz{Communication templates} para revisar tipo, estado activo y errores de validacion.}
\paso{2}{Abra \campointerfaz{Communication recipients} para buscar destinatarios por nombre, correo o institucion.}
\paso{3}{Abra \campointerfaz{Communication batches} para revisar tipo, estado, modo, correo de prueba, creador y fechas.}
\paso{4}{Abra \campointerfaz{Communication logs} para revisar destinatario, correo original, correo fisico, tipo, estado, error y fecha.}
\paso{5}{No modifique logs para ocultar errores. Los logs son trazabilidad operativa.}

\begin{resultadoesperado}
El superusuario puede reconstruir que plantilla, destinatarios y modo se usaron en una comunicacion sin ejecutar nuevos envios.
\end{resultadoesperado}

\section{Operaciones que requieren autorizacion}

\begin{tabularx}{\linewidth}{>{\bfseries}p{4.1cm}X}
\toprule
Operacion & Riesgo\\
\midrule
Editar registros de participantes & Puede cambiar datos de contacto, categoria, asistencia o integrantes sin pasar por formularios publicos.\\
Ejecutar acciones masivas & Puede enviar solicitudes o sincronizar equipos de muchos registros a la vez.\\
Editar equipos o miembros & Puede alterar el cuadro oficial si la competencia ya fue generada.\\
Editar grupos, batallas o estados & Puede romper la coherencia entre resultados, participantes vivos e historial.\\
Editar plantillas o destinatarios & Puede afectar comunicaciones futuras o exponer datos personales.\\
Eliminar registros & Puede perder trazabilidad y relaciones. Evitelo salvo procedimiento tecnico autorizado.\\
\bottomrule
\end{tabularx}

\section{Buenas practicas}

\begin{itemize}
    \item Use el panel operativo para operar el torneo y Admin para soporte.
    \item Filtre antes de seleccionar registros.
    \item Lea mensajes administrativos despues de cada accion.
    \item No capture pantallas con correos, documentos o tokens reales.
    \item No cree duplicados para resolver errores de busqueda.
    \item Revise el panel operativo despues de cualquier correccion directa.
    \item Mantenga una bitacora externa de cambios administrativos sensibles.
\end{itemize}

\section{Checklist del capitulo}

\begin{itemize}
    \item Entorno correcto verificado antes de entrar.
    \item Modelo identificado antes de abrir o editar.
    \item Filtros aplicados antes de acciones masivas.
    \item Mensajes de accion leidos y registrados.
    \item Cambios directos evitados cuando existe flujo operativo.
    \item Comunicaciones auditadas sin modificar logs.
    \item Panel de organizacion revisado despues de correcciones administrativas.
\end{itemize}
